% ----------------------------------------------------------------------
%                   LATEX TEMPLATE FOR PhD THESIS
% ----------------------------------------------------------------------

% based on Harish Bhanderi's PhD/MPhil template, then Uni Cambridge
% http://www-h.eng.cam.ac.uk/help/tpl/textprocessing/ThesisStyle/
% corrected and extended in 2007 by Jakob Suckale, then MPI-CBG PhD programme
% and made available through OpenWetWare.org - the free biology wiki


%: Style file for Latex
% Most style definitions are in the external file PhDthesisPSnPDF.
% In this template package, it can be found in ./Latex/Classes/
\documentclass[twoside,11pt]{Latex/Classes/PhDthesisPSnPDF}


%: Macro file for Latex
% Macros help you summarise frequently repeated Latex commands.
% Here, they are placed in an external file /Latex/Macros/MacroFile1.tex
% An macro that you may use frequently is the figuremacro (see introduction.tex)
% This file contains macros that can be called up from connected TeX files
% It helps to summarise repeated code, e.g. figure insertion (see below).

% insert a centered figure with caption and description
% parameters 1:filename, 2:title, 3:description and label
\newcommand{\figuremacro}[3]{
	\begin{figure}[htbp]
		\centering
		\includegraphics[width=1\textwidth]{#1}
		\caption[#2]{\textbf{#2} - #3}
		\label{#1}
	\end{figure}
}

% insert a centered figure with caption and description AND WIDTH
% parameters 1:filename, 2:title, 3:description and label, 4: textwidth
% textwidth 1 means as text, 0.5 means half the width of the text
\newcommand{\figuremacroW}[4]{
	\begin{figure}[htbp]
		\centering
		\includegraphics[width=#4\textwidth]{#1}
		\caption[#2]{\textbf{#2} - #3}
		\label{#1}
	\end{figure}
}

% inserts a figure with wrapped around text; only suitable for NARROW figs
% o is for outside on a double paged document; others: l, r, i(inside)
% text and figure will each be half of the document width
% note: long captions often crash with adjacent content; take care
% in general: above 2 macro produce more reliable layout
\newcommand{\figuremacroN}[3]{
	\begin{wrapfigure}{o}{0.5\textwidth}
		\centering
		\includegraphics[width=0.48\textwidth]{#1}
		\caption[#2]{{\small\textbf{#2} - #3}}
		\label{#1}
	\end{wrapfigure}
}

% predefined commands by Harish
\newcommand{\PdfPsText}[2]{
  \ifpdf
     #1
  \else
     #2
  \fi
}

\newcommand{\IncludeGraphicsH}[3]{
  \PdfPsText{\includegraphics[height=#2]{#1}}{\includegraphics[bb = #3, height=#2]{#1}}
}

\newcommand{\IncludeGraphicsW}[3]{
  \PdfPsText{\includegraphics[width=#2]{#1}}{\includegraphics[bb = #3, width=#2]{#1}}
}

\newcommand{\InsertFig}[3]{
  \begin{figure}[!htbp]
    \begin{center}
      \leavevmode
      #1
      \caption{#2}
      \label{#3}
    \end{center}
  \end{figure}
}


%%% Local Variables: 
%%% mode: latex
%%% TeX-master: "~/Documents/LaTeX/CUEDThesisPSnPDF/thesis"
%%% End: 


%: ----------------------------------------------------------------------
%:                  TITLE PAGE: name, degree,..
% ----------------------------------------------------------------------
% below is to generate the title page with crest and author name

%if output to PDF then put the following in PDF header
\ifpdf  
    \pdfinfo { /Title  (Dise\~{n}o y Fabricaci\'{o}n De Un Pico-Sat\'{e}lite Artificial Mediante Las Normas CanSat)
               /Creator (TeX)
               /Producer (pdfTeX)
               /Author (Israel Zu\~{n}iga de la Mora israelzuniga@ucol.mx)
               /CreationDate (D:YYYYMMDDhhmmss)  %format D:YYYYMMDDhhmmss
               /ModDate (D:YYYYMMDDhhmm)
               /Subject (xyz)
               /Keywords (add, your, keywords, here) }
    \pdfcatalog { /PageMode (/UseOutlines)
                  /OpenAction (fitbh)  }
\fi


\title{Dise\~{n}o y Fabricaci\'{o}n De Un Pico-Sat\'{e}lite Artificial Mediante Las Normas CanSat}



% ----------------------------------------------------------------------
% The section below defines www links/email for author and institutions
% They will appear on the title page of the PDF and can be clicked
\ifpdf
  \author{\href{mailto:israelzuniga@ucol.mx}{Israel Zu\~{n}iga de la Mora}}
%  \cityofbirth{born in XYZ} % uncomment this if your university requires this
%  % If city of birth is required, also uncomment 2 sections in PhDthesisPSnPDF
%  % Just search for the "city" and you'll find them.
  \collegeordept{\href{http://telematicanet.ucol.mx}{Facultad de Telem\'{a}tica}}
  \university{\href{http://www.ucol.mx}{Universidad de Colima}}

  % The crest is a graphics file of the logo of your research institution.
  % Place it in ./0_frontmatter/figures and specify the width
  \crest{\includegraphics[width=4cm]{escudo-ucol}}
  
% If you are not creating a PDF then use the following. The default is PDF.
\else
  \author{Israel Zu\~{n}iga de la Mora}
%  \cityofbirth{born in XYZ}
  \collegeordept{Facultad de Telem\'{a}tica}
  \university{Universidad de Colima}
  \crest{\includegraphics[width=4cm]{escudo-ucol}}
\fi

%\renewcommand{\submittedtext}{change the default text here if needed}
\degree{Ingenier\'{i}a en Telem\'{a}tica}
\degreedate{Junio 2011}


% ----------------------------------------------------------------------
       
% turn of those nasty overfull and underfull hboxes
\hbadness=10000
\hfuzz=50pt


%: --------------------------------------------------------------
%:                  FRONT MATTER: dedications, abstract,..
% --------------------------------------------------------------

\begin{document}

%\language{english}

% sets line spacing
\renewcommand\baselinestretch{1.2}
\baselineskip=18pt plus1pt


%: ----------------------- generate cover page ------------------------

\maketitle  % command to print the title page with above variables


%: ----------------------- cover page back side ------------------------
% Your research institution may require reviewer names, etc.
% This cover back side is required by Dresden Med Fac; uncomment if needed.

%\newpage
\vspace{10mm}
1. Asesor: Profesor Omar \'{A}lvarez C\'{a}rdenas

%\vspace{10mm}
%2. Reviewer: 

%\vspace{20mm}
%Day of the defense:

%\vspace{20mm}
%\hspace{70mm}Signature from head of PhD committee:



%: ----------------------- abstract ------------------------

% Your institution may have specific regulations if you need an abstract and where it is to be placed in the document. The default here is just after title.

%martes 7 de junio - suprimido abstract
%
% Thesis Abstract -----------------------------------------------------


%\begin{abstractslong}    %uncommenting this line, gives a different abstract heading
\begin{abstracts}        %this creates the heading for the abstract page

Put your abstract or summary here, if your university requires it.

\end{abstracts}
%\end{abstractlongs}


% ---------------------------------------------------------------------- 


% The original template provides and abstractseparate environment, if your institution requires them to be separate. I think it's easier to print the abstract from the complete thesis by restricting printing to the relevant page.
% \begin{abstractseparate}
%   
% Thesis Abstract -----------------------------------------------------


%\begin{abstractslong}    %uncommenting this line, gives a different abstract heading
\begin{abstracts}        %this creates the heading for the abstract page

Put your abstract or summary here, if your university requires it.

\end{abstracts}
%\end{abstractlongs}


% ---------------------------------------------------------------------- 

% \end{abstractseparate}


%: ----------------------- tie in front matter ------------------------

%martes 7 de junio - suprimidos

\frontmatter
\include{}
\include{}


%: ----------------------- contents ------------------------

\setcounter{secnumdepth}{3} % organisational level that receives a numbers
\setcounter{tocdepth}{3}    % print table of contents for level 3
\tableofcontents            % print the table of contents
% levels are: 0 - chapter, 1 - section, 2 - subsection, 3 - subsection


%: ----------------------- list of figures/tables ------------------------

\listoffigures	% print list of figures

\listoftables  % print list of tables


%: ----------------------- glossary ------------------------

% Tie in external source file for definitions: /0_frontmatter/glossary.tex
% Glossary entries can also be defined in the main text. See glossary.tex
% this file is called up by thesis.tex
% content in this file will be fed into the main document

% Glossary entries are defined with the command \nomenclature{1}{2}
% 1 = Entry name, e.g. abbreviation; 2 = Explanation
% You can place all explanations in this separate file or declare them in the middle of the text. Either way they will be collected in the glossary.

% required to print nomenclature name to page header
\markboth{\MakeUppercase{\nomname}}{\MakeUppercase{\nomname}}


% ----------------------- contents from here ------------------------

% chemicals
\nomenclature{DAPI}{4',6-diamidino-2-phenylindole; a fluorescent stain that binds strongly to DNA and serves to marks the nucleus in fluorescence microscopy} 
\nomenclature{DEPC}{diethyl-pyro-carbonate; used to remove RNA-degrading enzymes (RNAases) from water and laboratory utensils}
\nomenclature{DMSO}{dimethyl sulfoxide; organic solvent, readily passes through skin, cryoprotectant in cell culture}
\nomenclature{EDTA}{Ethylene-diamine-tetraacetic acid; a chelating (two-pronged) molecule used to sequester most divalent (or trivalent) metal ions, such as calcium (Ca$^{2+}$) and magnesium (Mg$^{2+}$), copper (Cu$^{2+}$), or iron (Fe$^{2+}$ / Fe$^{3+}$)}



 

\begin{multicols}{2} % \begin{multicols}{#columns}[header text][space]
\begin{footnotesize} % scriptsize(7) < footnotesize(8) < small (9) < normal (10)

%\printnomenclature[1.5cm] % [] = distance between entry and description
\label{nom} % target name for links to glossary

\end{footnotesize}
\end{multicols}



%: --------------------------------------------------------------
%:                  MAIN DOCUMENT SECTION
% --------------------------------------------------------------

% the main text starts here with the introduction, 1st chapter,...
\mainmatter

%\renewcommand{\chaptername}{} % uncomment to print only "1" not "Chapter 1"
\renewcommand{\chaptername}{Cap\'{i}tulo}


%: ----------------------- subdocuments ------------------------

% Parts of the thesis are included below. Rename the files as required.
% But take care that the paths match. You can also change the order of appearance by moving the include commands.


% this file is called up by thesis.tex
% content in this file will be fed into the main document

%: ----------------------- introduction file header -----------------------
\chapter{Introducci\'{o}n}

% the code below specifies where the figures are stored
\ifpdf
    \graphicspath{{1_introduction/figures/PNG/}{1_introduction/figures/PDF/}{1_introduction/figures/}}
\else
    \graphicspath{{1_introduction/figures/EPS/}{1_introduction/figures/}}
\fi

% ----------------------------------------------------------------------
%: ----------------------- introduction content ----------------------- 
% ----------------------------------------------------------------------



%: ----------------------- HELP: latex document organisation
% the commands below help you to subdivide and organise your thesis
%    \chapter{}       = level 1, top level
%    \section{}       = level 2
%    \subsection{}    = level 3
%    \subsubsection{} = level 4
% note that everything after the percentage sign is hidden from output



\section{CanSat} % section headings are printed smaller than chapter names
% intro
CanSat es el acr\'{o}nimo de “Canned Satellite” (Sat\'{e}lite Enlatado). El concepto fue introducido por primera vez a finales de la d\'{e}cada de 1990 en la Universidad de Stanford [2]. La idea detr\'{a}s del proyecto fue permitir a estudiantes enfrentarse a algunas situaciones encontradas en construir un sat\'{e}lite, pero al mismo tiempo para ser resueltas en un periodo de tiempo menor y con un gasto m\'{i}nimo. Los estudiantes deb\'{i}an dise\~{n}ar y construir instrumentos, adaptarlos a una lata de bebidas para despu\'{e}s lanzarla en un cohete. Entonces la lata de aluminio junto a su instrumentaci\'{o}n caen al suelo atados de un paraca\'{i}das mientras realizan diferentes tipos de experimentos. Todo esto siendo posible en menos de un a\~{n}o y a un costo muy bajo.


Cada a\~{n}o se celebra una competencia CanSat en Estados Unidos de Am\'{e}rica (“Annual CanSat Competition”, Texas). Esta competencia permite a equipos de diferentes universidades y preparatorias tener los retos de dise\~{n}ar y construir un sistema aeroespacial, de acuerdo a las especificaciones dadas por el comit\'{e} organizador. As\'{i} tambi\'{e}n se celebra otra competencia de interacci\'{o}n internacional, llamada ARLISS \'{o} “A Rocket Launch For International Student Satellites”.
As\'{i} mismo, toman lugar otras competencias CanSat en distintos pa\'{i}ses de Europa y As\'{i}a. Siendo el continente europeo [12] donde un mayor \'{i}ndice de estudiantes y pa\'{i}ses son involucrados.



\section{Normas CanSat} % subsection headings are again smaller than section names
Se tendr\'{a} un mejor entendimiento de \'{e}ste proyecto al revisar con calma las normas definidas para un concurso de pico-sat\'{e}lites CanSat. De manera resumida, las normas explican el valor educativo de una competencia CanSat, descripci\'{o}n de la misi\'{o}n, agrupaci\'{o}n de estudiantes para los equipos, diferentes tipos de misi\'{o}n CanSat, requerimientos, clases de CanSat, evaluaci\'{o}n y fases del proyecto.


\subsection{Requerimientos Generales}
En base a las gu\'{i}as de diferentes concursos se puede determinar:
\begin{enumerate}
\item La instrumentaci\'{o}n debe ser construida para caber dentro de una lata de bebida de 12 oz. (355ml). 115mm de altura y 66 mm de diametro.
\item La masa m\'{a}xima est\'{a} limitada a 350 gramos.
\item El costo m\'{a}ximo es de 1000 USD.
\item El paraca\'{i}das debe ser colocado en la parte superior y las antenas en el fondo del CanSat.
\item Cada CanSat debe contar con su propia estaci\'{o}n en tierra para fines de comunicaci\'{o}n.
\item El CanSat debe ser apagado antes del vuelo y durante su colocaci\'{o}n en \'{o}rbita. Solo debe ser activado despu\'{e}s de su separaci\'{o}n de el cohete o globo. [12]
\item El CanSat debe resistir una aceleraci\'{o}n de hasta 20G
\item Los materiales explosivos o de f\'{a}cil combusti\'{o}n dentro del dispositivo est\'{a}n prohibidos.
\end{enumerate}

\subsection{Clasificaci\'{o}n de Sat\'{e}lites por su tama\~{n}o}
El termino peque\~{n}o sat\'{e}lite es usado para mencionar a todo aquel s\'{a}telite con una masa menor a 500kg. A\'{u}n as\'{i} una  clasificaci\'{o}n m\'{a}s detallada es la propuesta:

%\begin{table}[htdp]
%\centering
%\begin{tabular}{ccc} % ccc means 3 columns, all centered; alternatives are l, r

%{\bf Gene} & {\bf GeneID} & {\bf Length} \\ 
% & denotes the end of a cell/column, \\ changes to next table row
%\hline % draws a line under the column headers

%human latexin & 1234 & 14.9 kbps \\
%mouse latexin & 2345 & 10.1 kbps \\
%rat latexin   & 3456 & 9.6 kbps \\
% Watch out. Every line must have 3 columns = 2x &. 
% Otherwise you will get an error.

%\end{tabular}
%\caption[title of table]{\textbf{title of table} - Overview of latexin genes.}
% You only need to write the title twice if you don't want it to appear in bold in the list of tables.
%\label{latexin_genes} % label for cross-links with \ref{latexin_genes}
%\end{table}
\begin{table}[Clasificaci\'{o}n de Peque\~{n}os Sat\'{e}lites]
\centering
\begin{tabular}{ | c | c |}
	\hline
	Clasificaci\'{o}n & Masa \\ \hline
	Mini-Sat\'{e}lite & 100-500kg \\
	Micro-Sat\'{e}lite & 10-100kg \\
	Nano-Sat\'{e}lite & 1-10kg \\ 
	Pico-Sat\'{e}lite & 0.1-1kg \\
	Femto-Sat\'{e}lite & $ < $ 100gr \\ \hline
\end{tabular}
\caption[Clasificaci\'{o}n de Peque\~{n}os Sat\'{e}lites]{\textbf{Clasficaci\'{o}n de Peque\~{n}os Sat\'{e}lites}}
\end{table}


\section{Planteamiento del Problema}
El desarrollo e investigaci\'{o}n en sistemas de software y hardware para sat\'{e}lites artificiales dentro de universidades es un tema acad\'{e}mico muy costoso; ya que usualmente se suelen planear proyectos para sat\'{e}lites de gran volumen, con tiempo de fabricaci\'{o}n a m\'{a}s de un a\~{n}o, un muy reducido tiempo de vida en el exterior y caracter\'{i}sticas de funcionamiento que  muchas veces son innecesarias para un proyecto realizado por estudiantes.


Por lo cual, en ocasiones este tipo de proyectos no puede verse funcionando en una \'{o}rbita al exterior del planeta. Principalmente por el costo total de fabricaci\'{o}n, lanzamiento al espacio mediante cohetes de transporte, colocaci\'{o}n en \'{o}rbita y notificaciones ante autoridades coordinadoras (‘Secretar\'{i}a de Comunicaciones y Transportes’, ‘Comisi\'{o}n Federal de Telecomunicaciones’, 'International Maritime Organization',  'International Civil Aviation Organization', 'Uni\'{o}n Internacional de Telecomunicaciones', ‘International Amateur Radio Union’ [4] y “Amateur Satellite Corporation’ [6] ). 


En M\'{e}xico los problemas existentes en las universidades y sector industrial para impulsar la investigaci\'{o}n en exploraci\'{o}n espacial y telecomunicaciones satelitales son: la escasez o inexistencia de presupuesto, tiempo, recursos humanos e intelectuales para estas \'{a}reas de conocimiento.


Otro aspecto importante es la calidad en la ense\~{n}anza y aprendizaje en las universidades p\'{u}blicas que ofrecen carreras de telecomunicaciones, el material de aprendizaje debe ser compartido entre los alumnos de un mismo grupo (antenas, transceptores, cables, conectores, decodificadores).

\section{Descripci\'{o}n de la Propuesta}
Un sistema de comunicaciones por sat\'{e}lite se divide en dos componentes: La estaci\'{o}n terrena y el sat\'{e}lite en \'{o}rbita. [10] La estaci\'{o}n terrena debe contar con al menos un transceptor de radio VHF y UHF, antenas para las bandas de radio frecuencia VHF y UHF, computadora personal para el procesamiento de telemetr\'{i}a y programaci\'{o}n de actividades para el sat\'{e}lite, conectores y adaptadores.


El segundo componente de la propuesta es el pico-sat\'{e}lite con un sistema b\'{a}sico, conocido como M\'{o}dulo de Servicio. Compuesto por los subsistemas de:

\begin{itemize}
\item Estructura
\item Telemetr\'{i}a
\item Alimentaci\'{o}n El\'{e}ctrica
\item Control Termal
\item Control de Posici\'{o}n
\end{itemize}

Y cumplir\'{a} con los requisitos generales para sat\'{e}lites CanSat que previamente han sido marcados.


La propuesta de trabajo consiste en la implementaci\'{o}n de los componentes, una estaci\'{o}n terrena y un pico-sat\'{e}lite CanSat.


Al concluir este trabajo de investigaci\'{o}n, la replicaci\'{o}n de la propuesta, por estudiantes de la Facultad de Telem\'{a}tica servir\'{a} para el aprendizaje y reafirmaci\'{o}n de conocimiento en las \'{a}reas de electr\'{o}nica, desarrollo de software, telecomunicaciones y direcci\'{o}n de proyectos.

\section{Preguntas de Investigaci\'{o}n}
\begin{itemize}
\item ¿Los resultados obtenidos ayudar\'{a}n a futuras investigaciones sobre peque\~{n}os sat\'{e}lites en el pa\'{i}s?
\item ¿El trabajo realizado podr\'{a} ser replicado por otros estudiantes de la Universidad de Colima?
\item ¿Los estudiantes que repliquen el proyecto se interesar\'{a}n por la radio afici\'{o}n?
\item ¿Qu\'{e} tecnolog\'{i}as implementar\'{a} este proyecto de investigaci\'{o}n?
\item ¿Qu\'{e} mejoras a tecnolog\'{i}as se podr\'{a}n generar al concluir de manera exitosa este proyecto de investigaci\'{o}n?
\item ¿Los costos de fabricaci\'{o}n y pruebas del CanSat ser\'{a}n costeables para incluirse como parte de los proyectos integradores en la Ingenier\'{i}a en Telem\'{a}tica?
\end{itemize}

\section{Objetivos}
\subsection{Objetivo General}
Implementar  un pico-sat\'{e}lite con el tama\~{n}o de una lata para bebidas de 12oz. con capacidades de env\'{i}o y recepci\'{o}n de informaci\'{o}n, o CanSat, para ser utilizado en experimentos cient\'{i}ficos, acad\'{e}micos y de prop\'{o}sito general.

\subsection{Objetivos Espec\'{i}ficos}
\begin{itemize}
\item Desarrollar un pico-sat\'{e}lite.
\item Colocar el CanSat en un vuelo sub-orbital a una distancia de al menos 20 mt de altura.
\item Documentar los est\'{a}ndares y procedimientos empleados para la fabricaci\'{o}n de sistemas aeroespaciales.
\item Adquirir informaci\'{o}n de telemetr\'{i}a mediante un enlace de radio a una estaci\'{o}n terrena una vez iniciado el viaje sub-orbital del prototipo.
\item Desarrollar materiales educativos que faciliten el conocimiento sobre exploraci\'{o}n espacial, sistemas de telecomunicaciones de uso amateur y comerciales basados en sat\'{e}lites artificiales, protecci\'{o}n y uso adecuado del entorno natural a partir de la fabricaci\'{o}n de sat\'{e}lites artificiales. 
\end{itemize}

\section{Hip\'{o}tesis}
Los resultados obtenidos con desarrollo del CanSat incrementaran el inter\'{e}s de los estudiantes de la Facultad de Telem\'{a}tica en la Universidad de Colima para continuar con el desarrollo e investigaci\'{o}n de sistemas de comunicaciones aeroespaciales.

\section{Justificaci\'{o}n}
Ante las dificultades dadas a conocer en el planteamiento del problema, la investigaci\'{o}n acad\'{e}mica puede verse beneficiada con la fabricaci\'{o}n y uso de peque\~{n}os sat\'{e}lites (clasificados como Nano, Pico y Femto sat.), siendo posible la investigaci\'{o}n y experimentaci\'{o}n de sistemas espaciales para estudiantes. El resultado del trabajo de investigaci\'{o}n, puede ser colocado a alturas de 1-20 km  mediante globos meteorol\'{o}gicos (globos de Helio) o a alturas de 20-1,000 mt mediante cohetes de modelismo espacial para representar una misi\'{o}n espacial.  Y a su vez, el peque\~{n}o sat\'{e}lite debe interactuar con las condiciones extremas de clima al funcionar fuera del globo terr\'{a}queo. 
Se pueden obtener grandes trabajos con un presupuesto limitado y los retos t\'{e}cnicos de optimizar cada recurso del sat\'{e}lite.

\section{Factibilidad}
Para realizar de manera exitosa \'{e}ste proyecto debe tener el sustento de factibilidad econ\'{o}mica y t\'{e}cnica. La factibilidad econ\'{o}mica determina que no se exceder\'{a} un presupuesto fijo o que se necesite una suma muy grande de dinero para solventar los gastos por el trabajo. La factibilidad t\'{e}cnica determina las capacidades del investigador y los recursos f\'{i}sicos o digitales requeridos.

\subsection{Factibilidad Econ\'{o}mica}
Un proyecto o misi\'{o}n bajo las gu\'{i}as de Competencia Cansat provee la oportunidad de experimentar el dise\~{n}o y ciclo de vida de un sistema aeroespacial. Dichas normas, de manera general, establecen un costo no mayor a de 1000 USD. El prototipo tendr\'{a} un costo menor al l\'{i}mite por ser fabricado con componentes que previamente se ten\'{i}an y otros pueden comprarse con descuentos para estudiante.

\subsection{Factibilidad T\'{e}cnica}


Hoy en d\'{i}a se tiene alcance a herramientas necesarias para el dise\~{n}o y fabricaci\'{o}n de dispositivos electr\'{o}nicos dentro de la Universidad de Colima. El acceso a Internet y el comercio electr\'{o}nico permiten la b\'{u}squeda y compra de componentes o materiales electr\'{o}nicos y electricos que agoten su existencia en las ciudades de Colima o Guadalajara; por lo cual no existe riesgo por una investigaci\'{o}n inconclusa debido a material faltante. El conocimiento de las herramientas de dise\~{n}o y fabricaci\'{o}n ha sido aprendido en cursos dentro de esta facultad.


%: ----------------------- HELP: special characters
% above you can see how special characters are coded; e.g. $\alpha$
% below are the most frequently used codes:
%$\alpha$  $\beta$  $\gamma$  $\delta$

%$^{chars to be superscripted}$  OR $^x$ (for a single character)
%$_{chars to be suberscripted}$  OR $_x$

%>  $>$  greater,  <  $<$  less
%≥  $\ge$  greater than or equal, ≤  $\ge$  lesser than or equal
%~  $\sim$  similar to

%$^{\circ}$C   ° as in degree C
%±  \pm     plus/minus sign

%$\AA$     produces  Å (Angstrom)




% dextran, starch, glycogen continued


%: ----------------------- HELP: references
% References can be links to figures, tables, sections, or references.
% For figures, tables, and text you define the target of the link with \label{XYZ}. Then you call cross-link with the command \ref{XYZ}, as above
% Citations are bound in a very similar way with \cite{XYZ}. You store your references in a BibTex file with a programme like BibDesk.




%: ----------------------- HELP: adding figures with macros
% This template provides a very convenient way to add figures with minimal code.
% \figuremacro{1}{2}{3}{4} calls up a series of commands formating your image.
% 1 = name of the file without extension; PNG, JPEG is ok; GIF doesn't work
% 2 = title of the figure AND the name of the label for cross-linking
% 3 = caption text for the figure

%: ----------------------- HELP: www links
% You can also see above how, www links are placed
% \href{http://www.something.net}{link text}

% variation of the above macro with a width setting
% \figuremacroW{1}{2}{3}{4}
% 1-3 as above
% 4 = size relative to text width which is 1; use this to reduce figures





%: ----------------------- HELP: lists
% This is how you generate lists in LaTeX.
% If you replace {itemize} by {enumerate} you get a numbered list.


 


%: ----------------------- HELP: tables
% Directly coding tables in latex is tiresome. See below.
% I would recommend using a converter macro that allows you to make the table in Excel and convert them into latex code which you can then paste into your doc.
% This is the link: http://www.softpedia.com/get/Office-tools/Other-Office-Tools/Excel2Latex.shtml
% It's a Excel template file containing a macro for the conversion.

%\begin{table}[htdp]
%\centering
%\begin{tabular}{ccc} % ccc means 3 columns, all centered; alternatives are l, r

%{\bf Gene} & {\bf GeneID} & {\bf Length} \\ 
% & denotes the end of a cell/column, \\ changes to next table row
%\hline % draws a line under the column headers

%human latexin & 1234 & 14.9 kbps \\
%mouse latexin & 2345 & 10.1 kbps \\
%rat latexin   & 3456 & 9.6 kbps \\
% Watch out. Every line must have 3 columns = 2x &. 
% Otherwise you will get an error.

%\end{tabular}
%\caption[title of table]{\textbf{title of table} - Overview of latexin genes.}
% You only need to write the title twice if you don't want it to appear in bold in the list of tables.
%\label{latexin_genes} % label for cross-links with \ref{latexin_genes}
%\end{table}



% There you go. You already know the most important things.


% ----------------------------------------------------------------------



	% background information
% this file is called up by thesis.tex
% content in this file will be fed into the main document

\chapter{Estado del Arte} % top level followed by section, subsection


% ----------------------- contents from here ------------------------

Los peque\~{n}os sat\'{e}lites han capturado la imaginaci\'{o}n a nivel mundial de fabricantes y dise\~{n}adores de los mismos, debido a sus capacidades econ\'{o}micas y tecnol\'{o}gicas. En las universidades con intereses en tecnolog\'{i}a espacial el primer foco de atenci\'{o}n son los objetos espaciales considerados como micro, nano y pico sat\'{e}lites. Seg\'{u}n avances en la miniaturizaci\'{o}n electr\'{o}nica se permite tener cargas \'{u}tiles (payloads) e instrumentos para misiones espaciales de menor tama\~{n}o y masa, tambi\'{e}n de un consumo de energ\'{i}a m\'{a}s eficiente. Siendo popular el est\'{a}ndar CubeSat [1] (acr\'{o}nimo de Cube Satellite) dise\~{n}ado principalmente por Jordi Puig-Suari y Robert Twiggs, adem\'{a}s creador del concepto CanSat.


Recientemente en el pa\'{i}s se cuenta con la referencia escasa de proyectos relacionados con peque\~{n}os sat\'{e}lites. Siendo el primero, un trabajo sin publicaciones o avances por parte del Instituto de Rob\'{o}tica de Yucat\'{a}n con el desarrollo de un CubeSat (TRIY-SAT I), en segunda instancia se puede relacionar al proyecto SARSEM-ICARUS II o Sistema Aerost\'{a}tico de Repetici\'{o}n Sub-Espacial Mexicano que realiz\'{o} un viaje sub-orbital mediante un globo meteorol\'{o}gico, desarrollado por el Club de Radio Amateur del Estado de Guanajuato A.C. y liderado por los ingenieros Erick Arzola (XE1CKJ) y Jonathan Remba (XE1BRX); con participaci\'{o}n y apoyo de operadores de radio amateur dentro del pa\'{i}s y el extranjero. Durante 2008 se realizan dos proyectos por estudiantes del Tecnol\'{o}gico de Monterrey, campus Puebla y campus Ciudad de M\'{e}xico respectivamente; estas dos iniciativas crearon un veh\'{i}culo espacial de acuerdo a las normas de competencia CanSat. Teniendo lugar en la “Annual Cansat Competition” de 2008 en Texas, EUA.





% ---------------------------------------------------------------------------
% ----------------------- end of thesis sub-document ------------------------
% ---------------------------------------------------------------------------						% estado del arte

% this file is called up by thesis.tex
% content in this file will be fed into the main document

\chapter{Metodolog\'{i}a} % top level followed by section, subsection



% the code below specifies where the figures are stored
\ifpdf
    \graphicspath{{3/figures/PNG/}{3/figures/PDF/}{3/figures/}}
\else
    \graphicspath{{3/figures/EPS/}{3/figures/}}
\fi

% ----------------------- contents from here ------------------------

\section{Desarrollo Orientado a Prototipos}

El desarrollo orientado a prototipos es un m\'{e}todo de investigaci\'{o}n que normalmente se usa para el desarrollo de Software, pero tambi\'{e}n puede ser usado para crear prototipos de hardware. Un prototipo es un modelo del comportamiento del sistema que puede ser usado para entender completamente o ciertos aspectos de \'{e}l y as\'{i} clasificar los requerimientos. Un prototipo es una representaci\'{o}n de un sistema, aunque no es un sistema completo, posee las caracter\'{i}sticas del sistema final o parte de ellas.

Las fases que comprenden el m\'{e}todo de desarrollo orientado a prototipos propuesta es:

\figuremacroW{metodo-desarrollo}{M\'{e}todo de Desarrollo propuesto}{Desarrollo Orientado a Prototipos}{}






















% ---------------------------------------------------------------------------
% ----------------------- end of thesis sub-document ------------------------
% ---------------------------------------------------------------------------			
% this file is called up by thesis.tex
% content in this file will be fed into the main document

\chapter{Calendario} % top level followed by section, subsection

% the code below specifies where the figures are stored
\ifpdf
    \graphicspath{{4/figures/PNG/}{4/figures/PDF/}{4/figures/}}
\else
    \graphicspath{{4/figures/EPS/}{4/figures/}}
\fi
% ----------------------- contents from here ------------------------

\section{Calendario de Actividades}

La programaci\'{o}n de actividades se refleja en el calendario o cronograma de actividad. Siendo el \'{u}nico integrante de las actividades el mismo autor de este proyecto. El \'{a}rea sombreada de las tablas determina la duraci\'{o}n estimada de la actividad, teniendo una X si fue realizada, o el campo vaci\'{o} si no se completo. Los campos que contienen un asterisco ( * ), son aquellas actividades programadas para  ser realizadas m\'{a}s adelante.
Destaca la revisi\'{o}n y correci\'{o}n del proyecto, que por atenci\'{o}n a otros pendientes escolares no se pudo realizar conforme lo acordado con el asesor. Esto afectando de gran manera al avance del trabajo.

\figuremacroW{1}{Calendario de Actividades realizadas y por terminar para Marzo-Junio}{}{}

En las actividades programadas para el verano solo han sido consideras las que requieren trabajo individual y no revisi\'{o}n por parte del asesor.

\figuremacroW{2}{Calendario de Actividades programadas para Julio-Septiembre}{}{}






% ---------------------------------------------------------------------------
% ----------------------- end of thesis sub-document ------------------------
% ---------------------------------------------------------------------------	
%\include{5/XYZ}
%\include{6/XYZ}

%% this file is called up by thesis.tex
% content in this file will be fed into the main document

\chapter{Discussion} % top level followed by section, subsection


% ----------------------- paths to graphics ------------------------

% change according to folder and file names
\ifpdf
    \graphicspath{{7/figures/PNG/}{7/figures/PDF/}{7/figures/}}
\else
    \graphicspath{{7/figures/EPS/}{7/figures/}}
\fi


% ----------------------- contents from here ------------------------






% ---------------------------------------------------------------------------
% ----------------------- end of thesis sub-document ------------------------
% ---------------------------------------------------------------------------               % discussion of results

%
% this file is called up by thesis.tex
% content in this file will be fed into the main document

\chapter{Materials \& methods} % top level followed by section, subsection


% ----------------------- paths to graphics ------------------------

% change according to folder and file names
\ifpdf
    \graphicspath{{8/figures/PNG/}{8/figures/PDF/}{8/figures/}}
\else
    \graphicspath{{8/figures/EPS/}{8/figures/}}
\fi

% ----------------------- contents from here ------------------------




 

% ---------------------------------------------------------------------------
%: ----------------------- end of thesis sub-document ------------------------
% ---------------------------------------------------------------------------



 






        % description of lab methods




% --------------------------------------------------------------
%:                  BACK MATTER: appendices, refs,..
% --------------------------------------------------------------

% the back matter: appendix and references close the thesis


%: ----------------------- bibliography ------------------------

% The section below defines how references are listed and formatted
% The default below is 2 columns, small font, complete author names.
% Entries are also linked back to the page number in the text and to external URL if provided in the BibTex file.

% PhDbiblio-url2 = names small caps, title bold & hyperlinked, link to page 
\begin{multicols}{2} % \begin{multicols}{ # columns}[ header text][ space]
\begin{tiny} % tiny(5) < scriptsize(7) < footnotesize(8) < small (9)

\bibliographystyle{Latex/Classes/PhDbiblio-url2} % Title is link if provided
%\renewcommand{\bibname}{Referencias} % changes the header; default: Bibliography

%\bibliography{9_backmatter/references.bib} % adjust this to fit your BibTex file


\end{tiny}
\end{multicols}

% --------------------------------------------------------------
% Various bibliography styles exit. Replace above style as desired.

% in-text refs: (1) (1; 2)
% ref list: alphabetical; author(s) in small caps; initials last name; page(s)
%\bibliographystyle{Latex/Classes/PhDbiblio-case} % title forced lower case
%\bibliographystyle{Latex/Classes/PhDbiblio-bold} % title as in bibtex but bold
%\bibliographystyle{Latex/Classes/PhDbiblio-url} % bold + www link if provided

%\bibliographystyle{Latex/Classes/jmb} % calls style file jmb.bst
% in-text refs: author (year) without brackets
% ref list: alphabetical; author(s) in normal font; last name, initials; page(s)

%\bibliographystyle{plainnat} % calls style file plainnat.bst
% in-text refs: author (year) without brackets
% (this works with package natbib)


% --------------------------------------------------------------

% according to Dresden med fac summary has to be at the end
%
% Thesis Abstract -----------------------------------------------------


%\begin{abstractslong}    %uncommenting this line, gives a different abstract heading
\begin{abstracts}        %this creates the heading for the abstract page

Put your abstract or summary here, if your university requires it.

\end{abstracts}
%\end{abstractlongs}


% ---------------------------------------------------------------------- 


%: Declaration of originality
%%martes 7 de junio - suprimido. no necesario aun


% Thesis statement of originality -------------------------------------

% Depending on the regulations of your faculty you may need a declaration like the one below. This specific one is from the medical faculty of the university of Dresden.

%\begin{declaration}        %this creates the heading for the declaration page

%I herewith declare that I have produced this paper without the prohibited assistance of third parties and without making use of aids other than those specified; notions taken over directly or indirectly from other sources have been identified as such. This paper has not previously been presented in identical or similar form to any other German or foreign examination board.

%The thesis work was conducted from XXX to YYY under the supervision of PI at ZZZ.

%\vspace{10mm}

%CITY,


%\end{declaration}


% ----------------------------------------------------------------------



\end{document}
